\section{Resource construction}

The pilot evidence base contains 55 purposively selected public discussions dated from 2022 through 2026. Selection sought conceptual coverage across platforms, workflow layers, incidents, design debates, documentation requests, product reports, vendor participation, and positive community outcomes. The thread served as the discovery unit. Fourteen deliberately difficult threads were segmented into 45 episodes so that interface access, deployment state, physical representation, observability, recovery, validation, and knowledge production could be represented separately when they occurred in one discussion.

The taxonomy contains ten constructs. B1 and B10 represent knowledge packaging and documentation or support. B2--B4 represent interface access, deployment identity, and physical resources. B5--B7 represent observability, recovery, and scheduling. B8 represents test--claim alignment. B9 represents AI context and physical feedback. Each code has inclusion, exclusion, primary-code, evidence, adjacency, and boundary rules.

All coding was performed by one analyst. The release reports no inter-rater statistic. A three-thread process-validation exercise identified ambiguities in the instrument and produced a primary-code tie-break, explicit read scopes, a formal episode unit, expected episode counts, and coherence tests. The answer-bearing key and a generated blind projection are released separately so that an independent coder can conduct a future pass without seeing the expected classifications.

Public-data handling follows a minimization strategy. The release contains derived records, short anonymized quotations, and source URLs. It excludes contributor handles, profile attributes, private messages, and a verbatim corpus. The source-audit ledger records approved quotation wording, context review, later-update checks, source role, and prohibited inference.
